\chapter{What Data Look Like \INDEV}

Sometimes you can spot data leakage by simply opening a file with an
ordinary utility, as we did in the last chapter when we used Apple's
Preview application to read geolocation data hidden inside a JPEG. 

Other times the leakage may be less obvious.  The software developers
know that the information is present but don't think that it's
relevant, and people who are in a position to know that the leakage is
relevant don't have the technical ability to determine that it's
present.

This chapter starts with instructions for creating a file that has hidden information
(\secref{sec:make-pdf}). We'll save the file in two ways, as a Rich
Text file (RTF) and as a Portable Document Format (PDF) file. Then
we'll look at the data the file contains using a variety of
techniques, including the raw bytes, a a \emph{hex dump}
(\secref{sec:hex-dumps}), and by extracting the printable strings
(\secref{sec:strings}). 

In order to understand what the hex dump and extracted strings
actually mean, we'll need to see how information such as 
(\secref{sec:numbers}) and letters (\secref{sec:letters}) are 
stored inside a computer's memory or in a file.  Much of this chapter is similar to
what might be covered at the start of a course on computer
architecture. But whereas an architecture course is primarily
concerned with computation and data in motion, this chapter is
concerned with understanding data at rest.

\section{Creating a File with ``Hidden'' text}\label{sec:make-pdf}

\figref{ch-what/textedit_hidden_1} shows how to create a simple file that
contains hidden data using Apple's TextEdit application.\footnote{The
  same set of steps will work with Libre Office or Microsoft Word. We
  use TextEdit here because the resulting RTF and PDF files are the
  simplest.}
 Start by
opening the TextEdit application, type the text, select the word ``Hidden'' and
change the background (highlight) color of the text from ``none'' to
``black.'' (The text background control is the box that looks like this:\includegraphics[height=1.5ex]{ch-what/textedit_highlight_control}).

Save two copies of the file, one in the Rich Text Format (RTF) in file
a called |hidden.rtf|, the second as an Adobe Portable Document Format
(PDF) file in a file named |hidden.pdf|. (You can use
the same approach with other word processors, but the resulting PDF
file is typically more complicated) 


\bifigure{ch-what/textedit_hidden_1}{ch-what/textedit_hidden_2}{Constructing a PDF with
  hidden text by selecting the text in a word processor (left) and
  then setting the highlight color to black (right).}

The file |hidden.rtf| is seven lines long and has just 297 bytes; that
file's contents appear in \lstref{hidden.rtf}. The file |hidden.pdf|, in
contrast, is 8304 bytes long. If the file's contents were printed here
they would look like a confusing jumble of letters, numbers, and
obscure symbols that would fill four pages. Such jumbled output is 
sometimes referred to as \emph{garbage}, and more properly termed
\emph{binary data} to distinguish it from the printable text that
makes up |hidden.rtf|. 

\sgraphic{ch-what/textedit_hidden-more}{The file {\tt hidden.pdf} viewed with the program \emph{less} running inside the Macintosh Terminal program.}

\lstinputlisting[caption=hidden.rtf,label=hidden.rtf]{ch-what/textedit_hidden.rtf}

\subsection{Viewing a text file}
\subsection{Viewing a binary file}

Attempt to view the file |hidden.pdf| with the
Unix program \emph{less} and you'll be warned:

\begin{code}
$ (@ more hidden.pdf @)
"hidden.pdf" may be a binary file. See it anyway?
\end{code}
%$

Type ``y'' and you'll see results similar to
\figref{ch-what/textedit_hidden-more}. The rest of this chapter explains what
this display means.

\section{Bits, Bytes, and Memory}\label{sec:numbers}
The \emph{bit} is the fundamental unit of storage on digital computers. Bits
have two states. Typically we call these |0| and |1|, but they can be any
two values (e.g. true/false, yes/no or up/down). This book 
uses the values |0| and |1| for bits. (The word \emph{bit} is a contraction of the words \emph{binary digit}.)

An ordered group of 4 bits is called a \emph{nibble}. Each bit can be
a |0| or a |1|, creating 16 possible different values. Note
that $2^4=16$. It is conventional to use the digits 0 through 9 to
number the first 10 possible nibble values, and use the letters A
through F to number the next six, as shown in \figref{nibble}. By
convention computer users call this notation \emph{hexadecimal},
although it's also called Base 16.

To avoid confusion between decimal and binary notation, this book will
always suffix binary numbers with a ``b''. Hexadecimal values have a
``h'' suffix. Some programming languages prefix hexadecimal strings
with a |\x| or |0x|. 

% \marginpar{There are 10 kinds of people in the world. Those who know binary, and those who don't.}


\begin{figure}
\begin{tabular}{>{\tt}c>{\tt}cc>{\tt}cl}
\textrm{binary} & \textrm{octal} & decimal & \textrm{hexadecimal} & english \\
\hline
0000b & 000 & 0 & 0h & zero \\
0001b & 001 & 1 & 1h & one \\
0010b & 002 & 2 & 2h & two \\
0011b & 003 & 3 & 3h & three \\
0100b & 004 & 4 & 4h & four \\
0101b & 005 & 5 & 5h & five \\
0110b & 006 & 6 & 6h & six \\
0111b & 007 & 7 & 7h & seven \\
1000b & 010 & 8 & 8h & eight \\
1001b & 011 & 9 & 9h & nine \\
1010b & 012 & 10& Ah & ten \\
1011b & 013 & 11& Bh & eleven \\
1100b & 014 & 12& Ch & twelve \\
1101b & 015 & 13& Dh & thirteen \\
1110b & 016 & 14& Eh & fourteen \\
1111b & 017 & 15& Fh & fifteen \\
\hline
\end{tabular}
\caption{A nibble is an ordered collection of 4 bits. Each bit can be
  a 0 or 1, so a nibble has $2^4=16$ possible values. This table shows
  the values of a nibble represented in binary, octal, decimal hexadecimal and English.}\label{nibble}
\end{figure}

Modern computers operate on 8 bits at a time, a unit called a
\emph{byte}. Because each bit in the byte can be a |0| or a |1|, and
because the bits are ordered, the byte can represent the values
\texttt{00000000b} through \texttt{11111111b} --- creating a
total of $2^8=256$ possible different values. This book will normally treat
bytes as unsigned integers in the range 0 to 255, but that is just a
convention; the Java programming language treats bytes as signed
integers in the range $-128$ to 127. It's also common to use present
the value of a byte in hexadecimal notation, with the range |00|
through |FF|.

\subsection{Registers}

Computers use groups of bytes to represent larger numbers. Two bytes
can represent 65,536 different values. Here are some examples:

\begin{tabular}{c>{\tt}c}
decimal & \textrm{binary} \\
\hline
0      & 00000000b 00000000b \\
255    & 00000000b 11111111b \\
256    & 00000001b 00000000b \\
65,280 & 11111111b 00000000b \\
65,535 & 11111111b 11111111b \\
\hline
\end{tabular}

Modern computers use registers in a variety of sizes. Many low-cost
embedded systems use 8-bit and 16-bit registers, most cell phones use
32-bit registers, and modern laptops, desktops and servers use
64-bit registers. There are also 80-bit registers used for some kinds
of scientific computing where both high resolution and high speed are required.

\subsection{Memory}

In the example above, the right-hand byte is called the \emph{least
  significant byte} and represents the bits $2^0$ through $2^7$; the
left-hand byte  called the \emph{most
significant byte}; it represents the binary values $2^8$ through
$2^{15}$. 

Because most computer memories are arranged as an array of bytes,
there are two ways to store these 16-bit values at a particular memory
location. 


% http://en.wikipedia.org/wiki/Single-precision_floating-point_format

Larger numbers Two bytes can be grouped together, allowing values between 0 and
65,535 to be represented. Four bytes

\section{Text, ASCII and Unicode}\label{sec:letters}

Modern computers do not work with bits directly, but use sequences of
8 bits called bytes. This sequence can represent
$2\times2\times2\times2\times2\times2\times2\times2=2^8=256$ distinct
values. The meaning assigned to these values generally depends on the
context. A few examples are shown in \tabref{tab:8bits}. The ASCII
codes represent NULL (no data is sent), SOH (Start Of Heading---the
beginning of a data header), and DEL (delete). NULL is a holdover from
the days of synchronous terminals, when there needed to be a constant
flow of data down all links: sending all 0s (NUL) meant that the data
should be ignored. DEL was chosen as delete when characters were
punched on paper tape and cards: in the event of an error, the
operator could backspace the punch and punch holes across the entire
line, causing the computer to ignore the value. We shall return
to the question of representation in \chapref{chap:unicode}.

\begin{table}
\begin{tabular}{cccccccc|c|c|c|c}
\multicolumn{8}{r}{                       Term:}  & unsigned char & signed char & ASCII \\
\multicolumn{8}{r}{                       Range:} & 0--255      & -128--127 & NULL--DEL\\
\cline{9-11}\multicolumn{8}{c}{Bits}                          &        &               &             &          \\
$2^7$ & $2^6$ & $2^5$ & $2^4$ & $2^3$ & $2^2$ & $2^1$ & $2^0$  &        &               &             &          \\
  128 &    64 &   32 &    16 &     8 &     4 &     2 &     1  &        &               &             &          \\
\hline
    0 &     0 &    0 &     0 &     0 &     0 &     0 &     0 &            0  &            0  &   NULL     \\
    0 &     0 &    0 &     0 &     0 &     0 &     0 &     1 &            1  &            1  &   SOH      \\
    0 &     1 &    0 &     0 &     0 &     0 &     0 &     1 &           65  &           65  &   A        \\
    0 &     1 &    1 &     0 &     0 &     0 &     0 &     1 &           96  &           65  &   a        \\
    0 &     1 &    1 &     1 &     1 &     1 &     1 &     1 &           127 &          127  &   DEL      \\
    1 &     0 &    0 &     0 &     0 &     0 &     0 &     0 &           128 &         -128  &   n/a      \\
    1 &     1 &    1 &     1 &     1 &     1 &     1 &     1 &           255 &           -1  &   n/a      \\
\end{tabular}
\caption{Eight bits can represent 256 distinct values.}\label{tab:8bits}
\end{table}



\section{Hex Dumps}\label{sec:hex-dumps}
A common way to visualize binary data is to format it as a series of
text lines, with each line containing a series of hexadecimal values (one
for each binary byte). Such displays are called \emph{hex
  dumps}.

There are many different programs that produce hex dumps, and each
program formats its output in a slight different manner. Fortunately
there is enough commonality that you can usually look at a hex dump
and figure it out. Most dumps will have a column on the left that
displays the offset in bytes from the beginning of the data, and a
column on the right that shows the ASCII representation of each
byte. Dumps are typically formatted such that each line represents 32
or 64 bytes of data.

Below is a hex dump for the 28-byte string 
|Hello World!\r\nHello World!\r\n| created with the Unix \emph{xxd}
program. Notice that the hexadecimal numbers \emph{do not} have a
``h'' suffix; users are expected to recognize hexadecimal values from
context:

\begin{Verbatim}
0000000: 4865 6c6c 6f20 576f 726c 6421 0d0a 4865  Hello World!..He
0000010: 6c6c 6f20 576f 726c 6421 0d0a            llo World!..
\end{Verbatim}


\section{Strings}\label{sec:strings}

Computer texts use the word \emph{string} inconsistently. Sometimes
the word refers to any sequence of zero or more bytes; sometimes it
means zero or more characters; and sometimes it means zero or
more printable characters. Usually the specific meaning can be
inferred from context. 

One quick way to analyze unfamiliar data is to look for printable
strings, which is the reason that ASCII is shown in most hex
dumps. With large blocks of data it can be hard to find the printable
strings. Consider the file random48.bin:

\begin{Verbatim}
0000000: fdf8 7c33 a153 0516 cc72 84fc 0050 7269  ..|3.S...r...Pri
0000010: 7661 6379 1437 ed4d e84b 0280 a74f e64b  vacy.7.M.K...O.K
0000020: 4695 6a66 9549 2cc9 4109 83c1 3e96 766b  F.jf.I,.A...>.vk
\end{Verbatim}

The file random48.bin contains the word ``Privacy'', but it can be hard to
find. The \emph{strings} program will extract printable
strings from a file and display them on stdout:

\begin{Verbatim}
$ strings random48.bin 
Privacy
$ 
\end{Verbatim}

Because ASCII is commonly used by computer systems, 
\emph{strings} command is a powerful way to find  identifying
information. 

\section{Magic Numbers}

\section{Encoding Binary Data}

BASE16, BASE64 and BASE85

\subsection{Legacy Encoding Methods}
uuencode

binhex


\section{High Entropy Data}
\subsection{Compressed Data}
\subsection{Encrypted Data}

\section{Hashing}

In computer science a \emph{hash function} is any function that maps a
sequence of zero or more characters (called a \emph{string}) to an
binary number of a specific fixed size---that is, a fixed number of
bits. A 16-bit hash function can produce $2^{16}=65,536$ different hash values, while a
32-bit hash function can produce $2^{32}=4,294,967,296$ possible hash
values. Hash functions are designed so that changing a single
character in the input results in a completely different number. While
the pigeonhole principle guarantees that many different strings will
have the same hash value---something that's called a \emph{hash
  collision}---the more bits in the hash has, the smaller the chance
of a collision.

Hashing was invented by H.\ P.\ Luhn in a 1953 IBM technical memo;
it's been widely used for computerized text processing since the
1960s. For example, because every sentence in a document can be
treated as a string, hashing makes it possible to rapidly see if the
same paragraph ever repeats in a long document: just computer the hash
value for each paragraph, put all of the hashes into a list, sort the
list, and see if any number repeats twice. If there is no repeat, then
no paragraph is repeated. Of course, if a number \emph{does} repeat,
then it's necessary to look at the paragraphs that correspond to those
two numbers to determine if the same paragraph really is repeated
twice or the duplicate is the result of a hash collision. Using hashes
in this manner is much faster than working directly with the
paragraphs because it is much faster for computers to compare numbers
than sequences of words---even when you take into account the time to
perform the hashing.

(Interesting fact: The name \emph{hash} comes from the way hash functions are typically
implemented as a two-step process that first chops and then mixes the
data, much in the way that one might make hash in the kitchen!)

In 1979 Stanford University PhD student Ralph Merkle invented a way to use
hashing for computer security\citep{merkle:79}. Merkle's idea was to
use a hash
function that produced more than 100 bits of output and that
additionally had the property of being \emph{one-way}. That is, it was
a function for which it was relatively easy to compute the hash of a
string, but it was nearly impossible, given a hash, to find a
corresponding string. The essence of Merkle's idea was to use a
document's 100-bit one-way hash as a stand-in for the document itself
for certain mathematical purposes. For example, instead of digitally
certifying a 50-page document, the document could be reduced to a
100-bit hash and the hash could then be certified. Because there are so many different possible hash values
($2^{100}\approx10^{30}$), Merkle reasoned that it would not be
possible for an attacker to take the digital signature from one document and use it
to certify a second document---because to do so would require that
both documents had the same hash value.

Today digital signatures applied to hashes are
the basis of many cyber security system. Digital signatures protect
credit card numbers sent over the Internet, certify the authenticity
and integrity of code run on iPhones, and validate keys used play
digital music. The idea of hashing has been applied to other areas as
well---in particular, forensics.

One of the first and continuing uses of hashing in DF was to establish
\emph{chain-of-custody} for forensic data. Instead of hashing a
document or a file, the hash function is applied to the entire disk
image---that is, \emph{all} of the bytes extracted from a hard
drive. Many law enforcement organizations will create two
disk images of a drive and then computer the hash of each image.  If the values match, then the
copies are assumed to each be a true copy of the data that were on the
drive. The hash value is then included in the official report. With
the hash recorded, the image file can be copied to a server at the
police department, to an investigator's workstation, or even given to
a forensic expert working for the other side. Any investigator with the
data can calculate the hash and see if it matches the original
reported value. If the hashes match, the investigator can be sure that
not a single bit in the disk image has been changed since the original
image was recorded. Hashing is so important that many DF tools can
automatically validate an evidence file by recomputing the hash and
comparing it with a stored value.

A second use for hashing is to identify specific
files. This approach takes advantage of the property that it is
extraordinarily unlikely for two files to have the same
hash value. File hashes can thus be used to identify files in much the
same way that a person can be identified by their fingerprints. 

Today forensic practitioners distribute databases containing the file
hashes as a standard part of forensic processing. The best known of these data
sets is the National Software Reference Library Reference Data Set,
distributed by the National Institute of Standards and
Technology\cite{nist-nsrl-rds-march2012}. These data sets can be used
to identify \emph{known goods}, such as programs distributed as part
of operating systems, or \emph{known bads} such as computer viruses,
stolen documents and child of pornography. Recent work is now
applying cryptographic hashing to blocks of data smaller than
files\cite{garfinkel:sector-id}, taking advantage of the fact that
even relatively short 512-byte and 4096-byte segments taken from
files can be highly identifying. 

File and sector identification with hashing means that a hard drive
containing millions of files can be automatically searched against a
database containing the hashes of hundreds of millions of file hashes
in a relatively short amount of time---perhaps just a few hours. Importantly, the search can be done without any human
intervention.

\section{Exercises}
* Convert numbers to different bases

* Look for private information with strings in files





